\documentclass[11pt]{article}
\usepackage[margin=1.08in]{geometry}
\usepackage{amsmath,amssymb,amsthm,mathtools}
\usepackage{booktabs,array,longtable}
\usepackage{xcolor}
\usepackage[colorlinks=true,linkcolor=blue!60!black,citecolor=blue!60!black,urlcolor=blue!60!black]{hyperref}

\newtheorem{theorem}{Theorem}[section]
\newtheorem{proposition}[theorem]{Proposition}
\newtheorem{lemma}[theorem]{Lemma}
\newtheorem{corollary}[theorem]{Corollary}
\newtheorem{cproposition}[theorem]{Computational Proposition}
\theoremstyle{definition}
\newtheorem{definition}[theorem]{Definition}
\newtheorem{conjecture}[theorem]{Field-Theoretic Hypothesis}
\theoremstyle{remark}
\newtheorem{remark}[theorem]{Remark}

\newcommand{\C}{\mathbb C}
\newcommand{\R}{\mathbb R}
\newcommand{\g}{\mathfrak g}
\newcommand{\W}{\mathcal W}
\newcommand{\F}{\mathcal F}
\newcommand{\cH}{\mathcal H}
\newcommand{\cL}{\mathcal L}
\newcommand{\cE}{\mathcal E}
\newcommand{\id}{\mathrm{id}}
\newcommand{\im}{\operatorname{im}}
\newcommand{\rank}{\operatorname{rank}}
\newcommand{\Tr}{\operatorname{Tr}}
\newcommand{\vol}{\operatorname{vol}}
\newcommand{\CS}{\operatorname{CS}}

\title{Residual $SO(4,2)$ Cohomology of Free Weyl Gravity\\
on the Conformal Cylinder:\\
Cartan Homotopy, Weyl-Square Classes, and Their Residual Pairing}
\author{Research-programme draft\\[-1mm]
\small Authorship and attribution to be finalized after external review}
\date{Draft for external review, July 2026}

\begin{document}
\maketitle

\begin{abstract}
Treating all fifteen cylinder conformal transformations, including compact
time translation $D$, as residual gauge transformations, we compute a
minimal absolute $SO(4,2)$ Chevalley--Eilenberg complex for free
four-dimensional Weyl gravity on $\R\times S^3$.  Its coefficient module is
the parity-complete on-shell Weyl-curvature Fock module.  The standard
invariant one-particle form is indefinite, with cylinder signs
$+I_E\oplus(-I_A)\oplus(-I_L)$.

The key mechanism is the Cartan homotopy formula.  The residual differential
obeys $d\iota_D+\iota_Dd=\cL_D$, so every subcomplex of nonzero total compact
degree is contractible.  In the Hamada-centered four-ghost polarization,
the selected absolute degree has ghost compact degree bounded below by
$-4$, and only matter weights at most four can contribute.  The vacuum and
one-particle sectors are acyclic, while a cutoff-complete absolute
calculation in the lowest two-particle sector gives
\[
 H^4_{\mathrm{res,min}}
 =\operatorname{span}\{[W_+^2],[W_-^2]\},
 \qquad G_{\rm res}=I_2.
\]
Here $G_{\rm res}=I_2$ is the induced Hermitian complementary-degree CE
pairing in the displayed residual basis.  The survivors are ghost-dressed vertex or deformation
classes, not one-particle graviton states.  Descent identifies their
parity-even and parity-odd combinations with
$C_{\mu\nu\rho\sigma}C^{\mu\nu\rho\sigma}$ and
$C_{\mu\nu\rho\sigma}\widetilde C^{\mu\nu\rho\sigma}$.  The first varies to
the Bach equation; the second is the Pontryagin boundary direction.  After
quotienting variationally trivial bulk functionals, the local dynamical
deformation space is one-dimensional and has residual Gram matrix $I_1$.

This algebraic residual theorem is exact once the on-shell Weyl module, the
Hamada ghost polarization, and full residual conformal gauging are assumed.
The last assumption is a physical choice: if $D$ is retained as a Hamiltonian
or boundary charge, the contraction is not a physical quotient.  The flat
adjoint BGG resolution and the topology of $\R\times S^3$ additionally give
an all-level smooth isomorphism between the Bach quotient and a geometrically
defined on-shell curvature system.  An exact symbolic cylinder construction
supplies finite-mode metric preimages for every $E/A/L$ block and identifies
that geometric system with the algebraic positive-energy module.  An
end-to-end polynomial calculation additionally constructs the metric,
ghost, equation-antifield and identity-antifield rows, measures the
noncompact homotopy defects, transfers the strict residual action, and
reproduces the centered cohomology without importing the hand-specified
$E/A/L$ matrices.  Two identifications remain conditional: the cross-energy
cyclic local-BV pairing and its normalization, and derivation of the centered
BFV polarization from the chosen closed-universe reduction.  Any analytic
completion requires additional control.  We claim neither a positive
graviton Fock space nor quantum anomaly freedom.
\end{abstract}

\tableofcontents

\section{Introduction}\label{sec:introduction}

Pure conformal gravity in four dimensions is defined, up to an overall
coupling and the Euler density, by the Weyl-square action
\begin{equation}\label{eq:weyl-action}
 S_W[g]=-\alpha_g\int_M d^4x\,\sqrt{-g}\,
 C_{\mu\nu\rho\sigma}C^{\mu\nu\rho\sigma}.
\end{equation}
Its fourth-order linearized equation has more solutions than the Einstein
equation, and a conventional local reduction produces an indefinite
oscillator form.  This familiar observation does not by itself determine the
state space after all gauge constraints have been imposed.  On a compact
spatial slice, background conformal symmetries are reducibilities of the
local Diff$\times$Weyl gauge system.  Treating those residual transformations
as gauge introduces a second, finite-dimensional reduction which is absent
from a naive oscillator count.  It is also a substantive physical choice.
In a Dirac or background-independence interpretation of a closed universe,
all residual conformal transformations can be quotiented.  In a fixed
cylinder conformal field theory, by contrast, $D$ is ordinarily the
Hamiltonian and the conformal group acts globally.  This paper computes the
first choice; it does not assert that strict pure-Weyl gravity uniquely
selects it.

The organizing question is therefore more basic than a sign assignment:
before asking whether a gravitational mode has positive or negative norm,
one must first determine whether it survives the full gauge reduction being
imposed.  The exact answer below concerns the residual reduction.  The
smooth metric-to-curvature bridge is proved globally, and an exact
all-energy cylinder construction identifies its algebraic positive-energy
part with the $E/A/L$ module used below.  The passage to the complete
local-plus-residual BV/BFV theory remains conditional.

The present paper isolates that free classical problem on
\begin{equation}
 M=\R\times S^3,
 \qquad
 \bar g=-dt^2+d\Omega_3^2,
\end{equation}
with unit cylinder radius.  Three spaces must be distinguished:
\begin{equation}\label{eq:three-levels}
 \begin{aligned}
 \text{local oscillator cohomology}
 &\supseteq \text{classically integrable perturbations},\\
 \text{classically integrable perturbations}
 &\longrightarrow \text{residual-BRST classes}.
 \end{aligned}
\end{equation}
The first contains the familiar Einstein, vector-descendant, and logarithmic
or higher-derivative towers.  The last is obtained only after the residual
$SO(4,2)$ constraints and their ghosts are included.  A linearized
oscillator need not survive this final reduction.

Our result is most naturally stated in the language of residual vertex
cohomology.  The Hamada ghost polarization supplies a top-degree four-ghost
factor of compact degree $-4$.  It pairs with scalar matter primaries of
weight four.  The complete centered residual cohomology consists of two
chiral Weyl-square classes, and their normalized residual CE Gram matrix is
positive.  Their parity combinations are the dynamical
Weyl-square density and the topological Pontryagin density.  Thus the
positive matrix is attached to classical vertex or theory-space directions.
It is not a positive norm on a propagating graviton Fock space; indeed, the
one-particle residual cohomology vanishes.

\paragraph{Relation to previous work.}
The conformal deformation detour complex on obstruction-flat backgrounds
was constructed by Branson and Gover \cite{BransonGover}; its formal
self-adjointness and related Yang--Mills detour constructions were developed
in \cite{GoverSombergSoucek}.  Gover and Peterson give the oriented
four-dimensional deformation sequence, while \v{C}ap proves that the flat
BGG sequence is a fine resolution and computes the same cohomology as the
flat twisted de Rham complex \cite{GoverPeterson,CapBGG}.  Cylinder
harmonics, oscillator charges, and
primary building blocks were developed by Hamada and Horata and by Hamada
\cite{HamadaHorata,HamadaBuildingBlocks}.  In the broader
conformal-mode/Riegert programme, Hamada identified the residual
diffeomorphisms with cylinder conformal transformations
\cite{HamadaCFT}, constructed the fifteen-generator residual BRST charge,
used the product of four proper-conformal ghosts as the centered ghost
vacuum, derived the scalar weight-four condition, and displayed a
Weyl-square BRST state \cite{HamadaBRST}.  None of those ingredients is
claimed as new here.  We use the harmonic and residual-algebra conventions,
but we do not import the Riegert field or the quantum claims of that model.
The Weyl-curvature module and its character are consistent with the general
conformal higher-spin description and partition-function literature
\cite{KuzenkoPonds,BeccariaBekaertTseytlin}.  Finally, the identification of
$C^2$ as the unique local four-derivative self-interaction of one linear
conformal graviton, under the assumptions stated there, is due to Boulanger
and Henneaux \cite{BoulangerHenneaux}.

The defensible new result is narrower and more concrete.  Starting with the
parity-complete pure-Weyl coefficient module, we compute the complete
minimal \emph{absolute} residual complex in the centered degree.  We prove
that the vacuum and one-particle sectors vanish, that precisely two
two-particle classes survive after the incoming-image test, and that their
normalized Hermitian residual CE Gram matrix is $I_2$.  The standard
Cartan homotopy makes the particle-number inventory exhaustive, and the
Euler--Lagrange map splits the two positive directions into one dynamical
and one topological class.

\paragraph{Claim boundary.}
Four statements must remain separate.
\begin{enumerate}
\item The \emph{algebraic residual theorem} is exact for the minimal
absolute $SO(4,2)$ Chevalley--Eilenberg complex built on the on-shell Weyl
module, the stated centered ghost polarization, and full residual conformal
gauging.
\item The smooth complexified metric-to-geometric-curvature bridge is an
exact global theorem.  Its restriction to the $D$-finite cylinder category
is also exact: the explicit $E/A/L$ harmonics give same-energy metric
preimages, nonzero curvature intertwiners, and character exhaustion.
Neither theorem is a statement about a completed Hilbert/Krein domain.
\item In the finite polynomial conformal category, an exact metric-BV
calculation now includes the ghost and antifield detour rows, constructs a
homotopy-equivariant retract, transfers the strict residual action, and
reproduces the centered kernel and image.  Identifying its representative
form with the complete cross-energy local-BV pairing remains conditional, as
does the residual BFV boundary choice which induces the centered ghost
insertion and its normalization.  These remaining hypotheses are displayed
rather than hidden.
\item Calling the resulting classes physical states additionally requires a
boundary and gauge principle which makes all fifteen conformal generators,
including $D$, gauge rather than global charges.  We call them
\emph{candidate physical classes under full residual conformal gauging}.
\end{enumerate}
No statement in this paper establishes quantum BRST nilpotency, anomaly
cancellation, interacting pseudo-unitarity, or a positive asymptotic
particle space.

The local-to-residual bridge belongs to the general subject of local BRST
cohomology with antifields, gauge fixing, and descent
\cite{BarnichBrandtHenneaux,BarnichHenneauxHurthSkenderis}.  Metsaev's
ordinary-derivative BRST--BV construction gives a concrete realization of
conformal fields and their conformal algebra on fields and antifields
\cite{MetsaevBRSTBV}.  Dneprov and Grigoriev construct a minimal nonlinear
Weyl-jet model and compatible presymplectic structure
\cite{DneprovGrigoriev}.  Their curvature-dependent ghost transformations
make that model a valuable comparison, but not an immediate proof that the
\emph{free} transferred residual differential is strict.  Neither work
supplies the global cylinder zero-mode normalization used here.

\section{The free local pure-Weyl complex}\label{sec:local-detour}

\subsection{Gauge, curvature, and equation maps}

Let $h_{\mu\nu}$ be a metric perturbation and let $(\xi^\mu,\sigma)$ be an
infinitesimal Diff$\times$Weyl parameter.  The full linear gauge map is
\begin{equation}\label{eq:full-K}
 K(\xi,\sigma)_{\mu\nu}
 =\cL_\xi\bar g_{\mu\nu}+2\sigma\bar g_{\mu\nu}
 =2\bar\nabla_{(\mu}\xi_{\nu)}+2\sigma\bar g_{\mu\nu}.
\end{equation}
After quotienting the trace, it becomes the conformal Killing operator
\begin{equation}\label{eq:K0}
 (K_0\xi)_{\mu\nu}
 =2\bar\nabla_{(\mu}\xi_{\nu)}
 -\frac12\bar g_{\mu\nu}\bar\nabla\!\cdot\!\xi.
\end{equation}
Let
\begin{equation}
 C_1h=\left.\frac{d}{d\varepsilon}
 C[\bar g+\varepsilon h]\right|_{\varepsilon=0}
\end{equation}
be the linearized Weyl-curvature map, and let $B_{\rm lin}$ denote the
action-normalized linearized Bach operator.  Formal adjoints use the
spacetime action pairing on a domain for which temporal and spatial boundary
terms vanish.

The conformal cylinder is conformally flat.  Naturality and conformal
covariance therefore give the exact local identity
\begin{equation}\label{eq:C1K}
 \boxed{C_1K=0}
\end{equation}
for arbitrary local gauge parameters, not only for the fifteen residual
conformal Killing parameters.

\begin{proposition}[Action-normalized factorization]\label{prop:factorization}
Use the repository normalization
\begin{equation}\label{eq:Sred}
 S_{\rm red}[g]=\int\sqrt{-g}
 \left(R_{\mu\nu}R^{\mu\nu}-\frac13R^2\right)
 =\frac12\int\sqrt{-g}\,(C^2-E_4).
\end{equation}
On a conformally flat background, modulo the Euler boundary term,
\begin{equation}\label{eq:Bfactor}
 \boxed{B_{\rm lin}=C_1^\sharp C_1.}
\end{equation}
Consequently,
\begin{equation}\label{eq:detour-identities}
 B_{\rm lin}K=0,
 \qquad
 K^\sharp B_{\rm lin}=0,
 \qquad
 B_{\rm lin}^\sharp=B_{\rm lin}.
\end{equation}
\end{proposition}

\begin{proof}
Since $C[\bar g]=0$, the mixed second variation of \eqref{eq:Sred} is
\begin{equation}
 \delta_h\delta_k S_{\rm red}
 =\langle C_1h,C_1k\rangle_{\W}.
\end{equation}
The definition
$\delta S_{\rm red}=\langle h,\cE(g)\rangle$ and
$B_{\rm lin}=D\cE|_{\bar g}$ gives
$\langle h,B_{\rm lin}k\rangle
=\langle C_1h,C_1k\rangle_{\W}$, proving
\eqref{eq:Bfactor}.  Equation \eqref{eq:C1K} then gives the first detour
identity, formal self-adjointness gives the second, and Hessian symmetry gives
the third.
\end{proof}

The Weyl-fibre pairing in Lorentz signature is indefinite.  Thus
$C_1^\sharp C_1$ is a formal-adjoint factorization, not a positivity
factorization.  In trace-free conformal geometry the corresponding local
sequence is
\begin{equation}\label{eq:detour-sequence}
 \Gamma(T)
 \xrightarrow{K_0}
 \Gamma(S^2_0T^*)
 \xrightarrow{B_{\rm lin}}
 \Gamma(S^2_0T^*)
 \xrightarrow{K_0^\sharp}
 \Gamma(T^*).
\end{equation}
On Bach-flat backgrounds this is the conformal deformation detour complex
of \cite{BransonGover}.  Formal self-adjointness holds in all signatures;
ellipticity is a Riemannian statement and is not assumed on the Lorentzian
cylinder.

\subsection{The global four-dimensional deformation complex}

The detour sequence \eqref{eq:detour-sequence} is not the only complex
available on the conformally flat cylinder.  On an oriented
four-dimensional conformally flat pseudo-Riemannian manifold, the adjoint
BGG deformation complex is \cite{GoverPeterson,CapBGG}
\begin{equation}\label{eq:BGG-sequence}
 \Gamma(T)
 \xrightarrow{K_0}
 \Gamma(S^2_0T^*[2])
 \xrightarrow{C_1}
 \Gamma(\mathcal C[2])
 \xrightarrow{C_1^\sharp\star}
 \Gamma(S^2_0T^*[-2])
 \xrightarrow{K_0^\sharp}
 \Gamma(T^*[-2]).
\end{equation}
Here $\mathcal C[2]$ is the algebraic Weyl-curvature bundle and
$C_1^\sharp\star$ means $C_1^\sharp\circ\star$, with the Hodge star acting
on one antisymmetric index pair.  In particular,
\begin{equation}\label{eq:star-complex-identity}
 \boxed{C_1^\sharp\star C_1=0.}
\end{equation}
The repository certificate verifies this identity directly as a
constant-coefficient differential-operator identity in both Euclidean and
Lorentzian conventions; Gover and Peterson supply the invariant global
sequence.

On a locally conformally flat manifold, the flat BGG complex is a fine
resolution of the sheaf of parallel sections of the flat adjoint-tractor
bundle \cite{CapDeformations,CapBGG}.  The cylinder
\begin{equation}
 M=\R\times S^3\simeq S^3
\end{equation}
is simply connected, so that flat local system is the constant system with
fibre $\mathfrak{so}(4,2)_\C$.  Fine-resolution cohomology therefore gives
an all-level global statement without an elliptic argument.

\begin{theorem}[Global smooth deformation cohomology]
\label{thm:global-BGG}
For smooth complexified global sections on the conformal cylinder,
\begin{equation}\label{eq:global-BGG-cohomology}
 H^q_{\rm def}(M)
 \cong H^q(S^3;\C)\otimes\mathfrak{so}(4,2)_\C
 \cong
 \begin{cases}
  \mathfrak{so}(4,2)_\C,&q=0,3,\\
  0,&q=1,2,4.
 \end{cases}
\end{equation}
In particular, the metric and curvature slots are globally exact:
\begin{equation}\label{eq:global-BGG-exactness}
 \boxed{\ker C_1=\im K_0,}
 \qquad
 \boxed{\ker(C_1^\sharp\star)=\im C_1.}
\end{equation}
\end{theorem}

\begin{proof}
The fine resolution computes sheaf cohomology of the flat adjoint local
system.  Simple connectivity trivializes its monodromy, while
$\R\times S^3$ deformation retracts onto $S^3$.  The only nonzero ordinary
cohomology groups are in degrees zero and three.  Vanishing in degrees one
and two is exactly the pair of statements
\eqref{eq:global-BGG-exactness}.
\end{proof}

The theorem is signature-independent and concerns unrestricted smooth
global sections.  It does not by itself establish continuity, closed range,
or bounded inverses on a Hilbert or Krein completion.

\subsection{The Bach quotient as an on-shell curvature space}

The passage from deformation-complex exactness to a detour equation is a
general exact-sequence argument.

\begin{proposition}[Detour--curvature lemma]\label{prop:detour-curvature}
Let
\begin{equation}
 E_0\xrightarrow{K}E_1\xrightarrow{C}E_2
 \xrightarrow{D_2}E_3
\end{equation}
be a complex which is exact at $E_1$ and $E_2$, and let $G:E_2\to E_1'$ be
an operator such that the detour equation is $B=GC$.  Then $[h]\mapsto Ch$
induces an isomorphism
\begin{equation}\label{eq:abstract-detour-curvature}
 \boxed{
 \frac{\ker B}{\im K}
 \cong \ker D_2\cap\ker G.}
\end{equation}
\end{proposition}

\begin{proof}
If $Bh=0$, then $GCh=0$, while $D_2Ch=0$ because the displayed sequence is
a complex.  If $Ch=0$, exactness at $E_1$ gives $h=K\xi$, so the induced map
is injective.  Conversely, if $U\in\ker D_2\cap\ker G$, exactness at $E_2$
gives $U=Ch$; then $Bh=GU=0$, proving surjectivity.
\end{proof}

The middle local cohomology relevant to the free field is
\begin{equation}\label{eq:local-middle}
 \mathcal H_B:=\frac{\ker B_{\rm lin}}{\im K_0}.
\end{equation}
In the four-dimensional deformation complex,
\begin{equation}
 D_2=C_1^\sharp\star,
 \qquad G=C_1^\sharp,
 \qquad B_{\rm lin}=GC_1.
\end{equation}
Define the geometrically on-shell curvature space
\begin{equation}\label{eq:Wgeom}
 \mathcal W_{\rm geom}
 :=\ker(C_1^\sharp\star)\cap\ker C_1^\sharp.
\end{equation}
The smooth global BGG theorem and
Proposition~\ref{prop:detour-curvature} identify the metric quotient with
this space exactly.

\begin{theorem}[Cylinder detour--curvature isomorphism]
\label{thm:curvature-isomorphism}
For smooth complexified fields on the conformal cylinder, the curvature map
$[h]\mapsto C_1h$ induces an $SO(4,2)$-equivariant isomorphism
\begin{equation}\label{eq:curvature-isomorphism}
 \boxed{
 \frac{\ker B_{\rm lin}}{\im K_0}
 \cong\mathcal W_{\rm geom}
 =\ker C_1^\sharp\cap\ker(C_1^\sharp\star).}
\end{equation}
\end{theorem}

\begin{proof}
If $B_{\rm lin}h=0$ and $U=C_1h$, then the factorization
\eqref{eq:Bfactor} gives $C_1^\sharp U=0$, while
\eqref{eq:star-complex-identity} gives $C_1^\sharp\star U=0$.  Thus the
curvature map lands in the displayed intersection.

If $C_1h=0$, exactness at the metric slot in
\eqref{eq:global-BGG-exactness} gives $h=K_0\xi$, proving injectivity.
Conversely, suppose $U$ lies in both kernels.  Exactness at the curvature
slot gives $U=C_1h$ for a smooth global $h$, and then
$B_{\rm lin}h=C_1^\sharp U=0$.  Hence every element of the intersection has
a Bach-flat metric potential, proving surjectivity.
\end{proof}

In Lorentz signature, complexified Weyl tensors split into Hodge
eigenspaces $U=U_++U_-$ with $\star U_\pm=\pm iU_\pm$.  If
$a=C_1^\sharp U_+$ and $b=C_1^\sharp U_-$, the two target equations are
\begin{equation}
 a+b=0,
 \qquad ia-ib=0,
\end{equation}
so $a=b=0$.  Define the smooth on-shell chiral geometric spaces by
\begin{equation}\label{eq:smooth-Weyl-spaces}
 \mathcal W_{{\rm geom},\pm}^{\rm sm}
 :=\{U_\pm:\star U_\pm=\pm iU_\pm,
                \ C_1^\sharp U_\pm=0\}.
\end{equation}
Then Theorem~\ref{thm:curvature-isomorphism} becomes
\begin{equation}\label{eq:smooth-local-target}
 \boxed{
 \frac{\ker B_{\rm lin}}{\im K_0}
 \cong
 \mathcal W_{{\rm geom},+}^{\rm sm}
 \oplus\mathcal W_{{\rm geom},-}^{\rm sm}.}
\end{equation}

All maps commute with $SO(4,2)$.  Hence the quotient-to-curvature map
preserves compact energy and both rotation spins on quotient classes.  It
remains to match this geometric system to the abstract positive-energy
module constructed in Section~\ref{sec:weyl-module}.

\begin{theorem}[Algebraic cylinder realization]\label{thm:algebraic-cylinder}
Let $\mathcal W_{{\rm geom},\pm}^{\rm alg,+}$ denote the $D$-finite,
$SO(4)$-finite positive-energy solutions of
\eqref{eq:smooth-Weyl-spaces}.  The curvature construction restricts to an
isomorphism
\begin{equation}\label{eq:local-target}
 \boxed{
 \mathcal W_{{\rm geom},+}^{\rm alg,+}
 \oplus\mathcal W_{{\rm geom},-}^{\rm alg,+}
 \cong\W_+\oplus\W_-.}
\end{equation}
Every curvature basis vector has an explicit Bach-flat metric representative
in the same $D$-energy and $SO(4)$ type.
\end{theorem}

\begin{proof}
Put $r=\tan(\beta/2)$, $z=1+r^2$, and
$q=d\beta+i\sin\beta\,d\gamma$.  For positive chirality the normalized
highest-weight metric representatives are the radiation-gauge cylinder
modes
\begin{equation}\label{eq:metric-preimage-families}
 \begin{array}{c|c|c}
  \text{tower}&J&\text{nonzero metric component}\\
  \hline
  E_n&n/2&h_{ij}=N_E(J)T[J]_{ij}e^{-int}\\
  A_n&(n-1)/2&h_{0i}=N_A(J)V[J]_ie^{-int}\\
  L_n&(n-2)/2&h_{ij}=N_L(J)T[J]_{ij}e^{-int}.
 \end{array}
\end{equation}
Here
\begin{align}
 V[J]&=-\frac{\sqrt{2J}}{4\pi}
 \cos^{2J-1}\!\frac\beta2\,
 e^{-i[(J+\frac12)\alpha+(J-\frac12)\gamma]}q,\nonumber\\
 T[J]&=\frac{\sqrt{2(2J-1)}}{16\pi}
 \cos^{2J-2}\!\frac\beta2\,
 e^{-i[(J+1)\alpha+(J-1)\gamma]}q\otimes q,
 \label{eq:metric-preimage-harmonics}
\end{align}
and
\begin{equation}
 N_E=\frac1{4\sqrt{J(2J+1)}},\qquad
 N_A=\frac1{2\sqrt{(2J-1)(2J+1)(2J+3)}},\qquad
 N_L=\frac1{4\sqrt{(J+1)(2J+1)}}.
\end{equation}

An exact coordinate calculation retains symbolic integer $n$, constructs
all components of $C_1h$, and applies $C_1^\sharp$ independently.  After
suppressing the common factor
$e^{-i(n t+m_L\alpha+m_R\gamma)}z^{-a}$, nonzero pivots are
\begin{align}
 (C_1h_E)_{0202}
 &=\frac{\sqrt{n(n-1)(n+1)}}{16\pi z^2},\nonumber\\
 (C_1h_A)_{0102}
 &=\frac{\sqrt{n(n-2)(n-1)}}{32\pi\sqrt{n+2}\,z},
 \label{eq:curvature-pivots}\\
 (C_1h_L)_{0202}
 &=\frac{\sqrt{n(n-3)(n-1)}}{16\pi z^2}.\nonumber
\end{align}
They are nonzero at every allowed tower energy.  The full tensors are
algebraically trace-free, have one common Hodge chirality, and obey
$C_1^\sharp C_1h=0$ identically in $n$.  Naturality of $C_1$ extends each
nonzero highest-weight intertwiner to its complete $SO(4)$ irrep.  Defining
$U_{F,n,M}=C_1h_{F,n,M}$ therefore gives the explicit same-block inverse
\begin{equation}\label{eq:cylinder-right-inverse}
 R_nU_{F,n,M}=h_{F,n,M},\qquad C_1R_n=\id,
 \qquad F\in\{E,A,L\}.
\end{equation}
The orientation-reversing cylinder isometry $\alpha\leftrightarrow\gamma$
exchanges the two $SU(2)$ factors and supplies the other chirality.

For exhaustion, the flat positive-chirality BGG data consist of the
lowest-weight curvature primary $(2;2,0)$, the equation module $(4;1,1)$,
and its identity $(5;\tfrac12,\tfrac12)$.  Their alternating character is
\begin{equation}
 \chi_{(2;2,0)}-\chi_{(4;1,1)}
 +\chi_{(5;\frac12,\frac12)}
 =\frac{5q^2-9q^4+4q^5}{(1-q)^4}.
\end{equation}
The cylinder module $\W_+$ is cyclic from the same primary and has this
character by \eqref{eq:weyl-character}; parity supplies the other chirality.
The nonzero graded intertwiner and equality of every graded $SO(4)$
multiplicity prove exhaustion.  Equations
\eqref{eq:metric-preimage-families}--\eqref{eq:cylinder-right-inverse}
supply finite harmonic metric representatives directly, excluding secular
or infinite-harmonic preimages in this algebraic category.
The exact coordinate tensor calculation, a machine-readable certificate,
and the generated pivot table are produced by
\path{symbolic/verify_conformal_cylinder_preimages.py}.
\end{proof}

\subsection{The finite-jet calculation as an independent convention audit}

The earlier finite-jet calculation remains useful, but it is no longer a
substitute for an unknown all-level theorem.  On homogeneous Euclidean
polynomial jets of degree $n=2,\ldots,6$, the matrices for $K$, $C_1$, and
$B_1$ obey
\begin{equation}
 C_1K=0,
 \qquad B_1K=0,
 \qquad \ker C_1=\im K.
\end{equation}
The exact ranks are shown in Table~\ref{tab:detour-ranks}.

\begin{table}[ht]
\centering
\small
\begin{tabular}{c@{\quad}rrrrrrrr}
\toprule
$n$ & $\dim G_n$ & $\dim h_n$ & $\dim W_{n-2}$ & $\dim B_{n-4}$
& $\rank K$ & $\rank C_1$ & $\rank B_1$
& $\dim(\ker B_1/\im K)$\\
\midrule
2 & 90  & 100 & 10  & 0   & 90  & 10  & 0  & 10\\
3 & 160 & 200 & 40  & 0   & 160 & 40  & 0  & 40\\
4 & 259 & 350 & 100 & 10  & 259 & 91  & 9  & 82\\
5 & 392 & 560 & 200 & 40  & 392 & 168 & 32 & 136\\
6 & 564 & 840 & 350 & 100 & 564 & 276 & 74 & 202\\
\bottomrule
\end{tabular}
\caption{Exact homogeneous-polynomial detour ranks.  The finite Euclidean
jet calculation independently audits the conventions and the first five
levels of the global BGG theorem.}
\label{tab:detour-ranks}
\end{table}

The quotient dimensions $(10,40,82,136,202)$ agree with the first five
levels of the parity-complete $E/A/L$ module introduced below.  The same
calculation separates a $15$-dimensional kernel of the low-degree gauge map,
generated by four translations, six rotations, one dilation, and four
special conformal transformations.  These finite results are independent
regression rails for \eqref{eq:global-BGG-exactness} and
\eqref{eq:local-target} in the repository normalization.

\subsection{The finite global zero-mode pair}

The same topology that kills the metric and curvature cohomology leaves two
finite-dimensional groups:
\begin{equation}\label{eq:BGG-zero-mode-pair}
 H^0_{\rm def}(M)\cong\mathfrak{so}(4,2)_\C,
 \qquad
 H^3_{\rm def}(M)\cong\mathfrak{so}(4,2)_\C.
\end{equation}
The degree-zero copy is the familiar space of fifteen conformal Killing
fields.  The degree-three copy lies on the adjoint equation side of the
self-dual deformation complex.  Refining both groups by compact weight gives
\begin{equation}
 4_{-1}\oplus7_0\oplus4_{+1}.
\end{equation}
Poincar\'e duality on $S^3$ and the Killing form pair the two copies, with
\begin{equation}
 \kappa(\g_r,\g_s)=0\qquad\text{unless}\qquad r+s=0.
\end{equation}

Its dimension, representation, and position are exactly those expected of
the dual residual constraint sector.  It is therefore a strong candidate
for the geometric source of the fifteen moment-map or Taub components
constructed in Section~\ref{sec:moment-map}.  That identification is not
yet a theorem: the full four-dimensional BGG sequence must first be kept
distinct from the shorter Bach detour complex, and the degree-three bundle
representative must be located explicitly.  Importantly, the direct
second-order argument in Proposition~\ref{prop:Taub-constraint} establishes
the Taub constraints without depending on this identification.  If the BGG
and BFV realizations are then compared, equivariance reduces the map to one
nonzero component and one overall scalar.

Thus the smooth deformation complex contains a controlled finite pair,
not an unspecified collection of extra global classes.  A complete BV
argument must still place the degree-three copy at the correct
ghost/antifield degree and construct the corresponding zero-mode projector.

\subsection{Domain boundary of the theorem}

Theorems~\ref{thm:global-BGG} and \ref{thm:curvature-isomorphism} are exact
for smooth complexified global sections.  Theorem~\ref{thm:algebraic-cylinder}
additionally identifies the $D$-finite, $SO(4)$-finite positive-energy
submodule by explicit metric representatives.  These theorems do not identify
spacelike-compact or distributional cohomology, nor do they prove
continuity and closed range on an infinite Hilbert/Krein completion.  If
such an analytic completion is required, Green-hyperbolic complexes and
causal homotopies remain natural tools \cite{GreenHyperbolic}.

The other remaining issue is BV rather than local geometry.  The fifteen
conformal-Killing zero modes and their degree-three duals must be split from
the local complex without double counting.  Dneprov and Grigoriev's minimal
nonlinear Weyl-jet model supplies compatible presymplectic data for this
comparison \cite{DneprovGrigoriev}, but its curvature-dependent ghost terms
must not be imported as a proof that the free residual transfer is strict.
It does not perform the global cylinder BFV split.

\section{The on-shell Weyl module}\label{sec:weyl-module}

\subsection{Character and cylinder towers}

We now define the coefficient module used in the exact residual
calculation.  A positive-chirality Weyl curvature is a conformal primary
with labels
\begin{equation}
 (\Delta;j_L,j_R)=(2;2,0),
\end{equation}
and the opposite chirality is its parity conjugate.  The Bach equation is in
$(4;1,1)$ and its differential identity is in
$(5;\tfrac12,\tfrac12)$.  The resulting one-chirality character is
\begin{equation}\label{eq:weyl-character}
 \chi_+(q)
 =\chi_{(2;2,0)}-\chi_{(4;1,1)}
 +\chi_{(5;\frac12,\frac12)}
 =\frac{5q^2-9q^4+4q^5}{(1-q)^4}.
\end{equation}
This defines the abstract algebraic positive-energy module $\W_+$.  The
character and tower calculation below, together with the explicit
intertwiner and metric preimages of
Theorem~\ref{thm:algebraic-cylinder}, identify it with the geometric chiral
system.

On $\R\times S^3$, $\W_+$ decomposes into three multiplicity-free towers:
\begin{align}
 E_n^+&:\quad
 \left(\frac{n+2}{2},\frac{n-2}{2}\right),
 &\dim E_n&=(n+3)(n-1), &&n\ge2,\label{eq:E-tower}\\
 A_n^+&:\quad
 \left(\frac n2,\frac{n-2}{2}\right),
 &\dim A_n&=(n+1)(n-1), &&n\ge3,\label{eq:A-tower}\\
 L_n^+&:\quad
 \left(\frac n2,\frac{n-4}{2}\right),
 &\dim L_n&=(n+1)(n-3), &&n\ge4.\label{eq:L-tower}
\end{align}
The negative-chirality towers exchange $j_L$ and $j_R$.  Summing
\eqref{eq:E-tower}--\eqref{eq:L-tower} reproduces
\eqref{eq:weyl-character}.  For both chiralities the first levels have
dimensions
\begin{equation}
 \dim\cH_2=10,
 \quad \dim\cH_3=40,
 \quad \dim\cH_4=82,
 \quad \dim\cH_5=136,
 \quad \dim\cH_6=202.
\end{equation}

The labels $E,A,L$ are cylinder oscillator labels only.  In particular,
they should not be interpreted as three independently signable conformal
representations.  Proper conformal transformations mix the towers.

\subsection{Invariant form and exact generator action}

In normalized cylinder coordinates the unique standard-real,
parity-compatible invariant form is, up to overall convention,
\begin{equation}\label{eq:Jconf}
 J_{\rm conf}=+I_E\oplus(-I_A)\oplus(-I_L).
\end{equation}
It is nondegenerate and indefinite.  The proper conformal lowering and
raising generators satisfy
\begin{equation}\label{eq:J-adjoint}
 (K^-_q)^\sharp=K^+_q,
 \qquad X^\sharp=X\quad(X=D,SU(2)_L,SU(2)_R),
\end{equation}
where $X^\sharp=J_{\rm conf}^{-1}X^\dagger J_{\rm conf}$.

Multiplicity one reduces each allowed lowering block to a scalar reduced
coefficient.  With source energy $n$, the six stable families are
\begin{align}
 c_{EE}(n)&=\sqrt{\frac{2(n-1)(n+1)(n+3)}{n+2}},
 &c_{AE}(n)&=\sqrt{\frac{8(n-1)}{(n-2)(n+2)}},\nonumber\\
 c_{AA}(n)&=\sqrt{\frac{2(n-3)(n-1)(n+2)}{n-2}},
 &c_{LE}(n)&=-\sqrt{\frac{2(n-3)}{n-2}},\label{eq:reduced-coefficients}\\
 c_{LA}(n)&=\sqrt{\frac8{n-2}},
 &c_{LL}(n)&=\sqrt{2(n-2)(n+1)}.\nonumber
\end{align}
Together with Clebsch--Gordan coefficients and \eqref{eq:J-adjoint}, these
determine the all-level one-particle $\mathfrak{so}(4,2)$ action.
The exact reconstruction verifies the conformal commutators on every
interior level of a finite energy buffer; the top shell is used only as a
buffer and is never mistaken for a finite conformal representation.

The free Fock coefficient module is the bosonic second quantization
\begin{equation}\label{eq:fock-module}
 \F_{\W}=\operatorname{Sym}(\W_+\oplus\W_-),
\end{equation}
with multiplicative lift of \eqref{eq:Jconf}.  The residual generators
preserve oscillator particle number.

\section{Hamiltonian moment-map normalization}\label{sec:moment-map}

The residual action is not introduced as an arbitrary matrix
representation.  In oscillator coordinates $z$, write the symplectic form
as
\begin{equation}\label{eq:symplectic-form}
 \Omega=i\,d\bar z\,J\wedge dz.
\end{equation}
If a conformal generator $X$ acts linearly by $\delta_Xz=K_Xz$, its
quadratic Hamiltonian moment map is
\begin{equation}\label{eq:moment-map}
 \mu_X(z,\bar z)=\bar z\,M_Xz,
 \qquad M_X=JK_X,
 \qquad d\mu_X=\iota_{X_X}\Omega.
\end{equation}
The $J$-adjoint identities ensure the correct reality of $\mu_X$.

\subsection{The conformal Taub constraint}

The same residual parameters arise as reducibilities of the local
Diff$\times$Weyl complex.  Let
\begin{equation}
 \epsilon=(\xi,\sigma),
 \qquad K(\xi,\sigma)=0,
 \qquad \sigma=-\frac14\bar\nabla_\mu\xi^\mu,
\end{equation}
and let $h$ satisfy $B^{(1)}[h]=0$.  With the Taylor convention
\begin{equation}
 g(\lambda)=\bar g+\lambda h+\frac{\lambda^2}{2}k+O(\lambda^3),
\end{equation}
the second-order equation is
\begin{equation}\label{eq:second-order-Bach}
 B^{(1)}[k]+B^{(2)}[h,h]=0.
\end{equation}

\begin{proposition}[Conformal Taub solvability constraint]
\label{prop:Taub-constraint}
For every residual conformal reducibility $\epsilon$, the current
\begin{equation}
 J_\epsilon^\mu[h]
 :=\xi_\nu B^{(2)\mu\nu}[h,h]
\end{equation}
is conserved.  If $h$ is tangent to a second-order family of Bach-flat
metrics on the compact cylinder, then
\begin{equation}\label{eq:Taub-charge}
 \boxed{
 \mathcal T_\epsilon[h]
 :=\int_{S^3}n_\mu\xi_\nu
 B^{(2)\mu\nu}[h,h],d\Sigma=0.}
\end{equation}
\end{proposition}

\begin{proof}
At second order around a Bach-flat background, the exact divergence and
trace identities imply
$\bar\nabla_\mu B^{(2)\mu\nu}[h,h]=0$ and
$B^{(2)\mu}{}_{\mu}[h,h]=0$ whenever $B^{(1)}[h]=0$.  Hence
\begin{equation}
 \bar\nabla_\mu J_\epsilon^\mu
 =B^{(2)\mu\nu}\bar\nabla_{(\mu}\xi_{\nu)}
 =-\sigma B^{(2)\mu}{}_{\mu}=0.
\end{equation}
If \eqref{eq:second-order-Bach} has a solution, the charge equals minus the
integral of $\xi_\nu B^{(1)\mu\nu}[k]$.  For a reducibility parameter, the
linearized Noether current is a spatial divergence.  Since
$\partial S^3=\varnothing$, its integral vanishes.  This is the standard
compact-Cauchy linearization-stability argument
\cite{AltasTekinLinearization,AltasTekinTaub}.
\end{proof}

The covariant phase-space Noether identity predicts that this same
quadratic constraint is the Hamiltonian of the residual linear action:
\begin{equation}\label{eq:moment-Taub-target}
 \mu_\epsilon(h)
 =\frac12\Omega_\Sigma(h,R_\epsilon h)
 \stackrel{?}{=}\pm\mathcal T_\epsilon[h].
\end{equation}
The sign and a possible factor of two depend only on the Taylor convention
for $B^{(2)}$.  Both sides form equivariant adjoint/coadjoint
$\mathfrak{so}(4,2)$ modules.  Since the complexified adjoint module is
simple, an equivariant nonzero comparison is unique up to one scalar.
Consequently the all-component equality can be proved by establishing
equivariance and matching one generator, preferably $D$.  This is also the
local-BRST relation between reducibilities and conserved-current classes
\cite{BarnichAntifield,BarnichBrandtHenneaux}.

In the curved-cylinder wave normalization used for the independent Taub
calculation, comparison with the quadratic action fixes the proper-conformal
kernel to
\begin{equation}\label{eq:taub-normalization}
 M_{\rm Taub}=-\frac{\sqrt2}{4\pi}\,J K^-
\end{equation}
up to the explicitly declared orientation and harmonic phases.  Two
independent curvature components fix the common coefficient; it is not a
one-entry fit.  An exact two-chirality jet through source energy four
verifies all fifteen moment-map components and the full conformal algebra on
interior energies two and three.

Equation \eqref{eq:taub-normalization} verifies
\eqref{eq:moment-Taub-target} for the tested proper-conformal kernels.  The
remaining all-level task is the single equivariant normalization comparison,
not fifteen unrelated curvature calculations.

This section supplies the Hamiltonian normalization of the residual action
and its Taub interpretation.  The cohomology theorem below depends only on
the exact module action, not on interpreting a finite moment-map jet as a
physical quotient.

\section{Residual BRST and Cartan homotopy reduction}\label{sec:cartan}

\subsection{The absolute residual complex}

Let
\begin{equation}
 \g=\mathfrak{so}(4,2)
\end{equation}
with compact grading
\begin{equation}\label{eq:g-grading}
 \g=\g_{-1}\oplus\g_0\oplus\g_{+1},
 \qquad \dim(\g_{-1},\g_0,\g_{+1})=(4,7,4).
\end{equation}
Here $D$ and $\mathfrak{so}(4)$ lie in $\g_0$, while proper conformal
lowering and raising generators lie in $\g_{-1}$ and $\g_{+1}$.
The dual ghosts carry the opposite compact degree.

For a coefficient module $V\subset\F_{\W}$, the absolute residual complex
is
\begin{equation}\label{eq:CE-complex}
 C^q(V)=V\otimes\Lambda^q\g^*,
\end{equation}
with Chevalley--Eilenberg differential
\begin{equation}\label{eq:CE-differential}
 d=c^a\rho(G_a)
 -\frac12 f^a{}_{bc}c^bc^c\iota_a.
\end{equation}
It raises ghost number by one and preserves the total compact degree
\begin{equation}\label{eq:total-degree}
 \delta=D_{\rm matter}+D_{\rm ghost}.
\end{equation}

The standard centered polarization uses the product $v$ of the four ghosts
dual to the raising generators.  It has
\begin{equation}\label{eq:ghost-vacuum-degree}
 \operatorname{gh}(v)=4,
 \qquad D_{\rm ghost}(v)=-4.
\end{equation}
This is the residual ghost convention used in the cylinder construction of
\cite{HamadaBRST}.  Its compact degree and uniqueness as the bottom
rotational scalar follow from the $(4,7,4)$ exterior-algebra polarization,
and Lemma~\ref{lem:ghost-pairing} shows that its complementary-degree pairing
is canonical up to one overall volume.  Only the match of that scalar to the
strict-pure-Weyl BV/BFV normalization remains a hypothesis in
Section~\ref{sec:transfer}; the broader Riegert sector of
\cite{HamadaBRST} is not included.

\subsection{Full residual gauging is a physical assumption}\label{sec:D-gauge}

Hamada's background-independence programme identifies the transformations
which preserve the cylinder gauge with residual diffeomorphisms and imposes
the resulting conformal charges through BRST
\cite{HamadaCFT,HamadaBRST}.  This supplies a direct precedent for the
complex \eqref{eq:CE-complex}.  It does not prove that every quantization of
strict pure-Weyl gravity must use the same quotient.

There are at least two distinct cylinder problems.
\begin{enumerate}
\item In the \emph{fully gauged residual problem}, the cylinder is a gauge
representative of a closed generally covariant system.  All fifteen
conformal-Killing transformations, including compact time translation, are
generated by residual constraints.  The ghost for $D$ exists and the Cartan
homotopy below is an allowed operation.
\item In the \emph{fixed-background cylinder problem}, $D$ is the
Hamiltonian and $SO(4,2)$ is a global symmetry acting on states or
observables.  Its eigenspaces are not quotiented.  The $D$ ghost and its
contraction cannot be used to erase nonzero energy.
\end{enumerate}
Proposition~\ref{prop:Taub-constraint} gives the first choice a direct
linearization-stability interpretation.  Within the closed-universe Dirac
boundary problem, the slice $S^3$ has no spatial boundary supporting an
independent asymptotic charge, and the residual conformal Killing parameters
instead generate quadratic integrability constraints, including the $D$
constraint.  Subject to the all-level normalization in
\eqref{eq:moment-Taub-target}, full residual gauging is therefore a coherent
zero-boundary-charge sector of pure conformal gravity; it is not a device
introduced merely to erase positive cylinder energy.  Compactness and the
Taub identity support this interpretation but do not make it the unique
physical quantization.

Adding a boundary or asymptotic region, coupling an external clock, or
deparametrizing the cylinder can convert $D$ into a physical global charge.
That is the second problem above, not a contradiction of the first.
Because $[K^-,K^+]$ contains $D$, the second problem is not obtained by
simply deleting one ghost while leaving the rest of
\eqref{eq:CE-differential} unchanged.  It requires a different relative or
boundary-charge complex.  Theorem~\ref{thm:residual-main} answers only the
first problem.  A strict pure-Weyl derivation must state the temporal
boundary conditions, the residual gauge group, and the treatment of its
charges before calling the result a physical state space.

\subsection{Cartan homotopy}

Let $\iota_D$ be contraction with the generator $D$ in the residual ghost
exterior algebra.  The usual Cartan identity is
\begin{equation}\label{eq:cartan-identity}
 d\iota_D+\iota_Dd=\cL_D.
\end{equation}

\begin{theorem}[Cartan homotopy reduction to compact degree zero]\label{thm:cartan}
Decompose the residual complex into total compact-degree eigenspaces,
\begin{equation}
 C=\bigoplus_{\delta}C_\delta,
 \qquad \cL_D|_{C_\delta}=\delta\,\id.
\end{equation}
Every $C_\delta$ with $\delta\ne0$ is contractible.  Hence inclusion of the
centered subcomplex is a quasi-isomorphism:
\begin{equation}\label{eq:cartan-quasiiso}
 \boxed{H^\bullet(C,d)\cong H^\bullet(C_{\delta=0},d).}
\end{equation}
\end{theorem}

\begin{proof}
On $C_\delta$ with $\delta\ne0$, define
\begin{equation}
 h_\delta=\frac{\iota_D}{\delta}.
\end{equation}
Equation \eqref{eq:cartan-identity} gives
\begin{equation}
 dh_\delta+h_\delta d=\id.
\end{equation}
Thus every closed vector of nonzero total degree has the explicit primitive
$\delta^{-1}\iota_D\Psi$.  If $\cL_D$ has a nilpotent part, the same proof
works wherever $\cL_D$ is invertible by using
$h=\iota_D\cL_D^{-1}$.
\end{proof}

The theorem depends on treating $D$ as a residual gauge generator.  If
cylinder time translation is retained as a physical Hamiltonian or a
boundary charge, contraction by its ghost is not an allowed quotient and
\eqref{eq:cartan-quasiiso} does not apply.

\subsection{Finite centered inventory}

At ghost number four, the minimal residual exterior algebra contains at
most the four negative-degree ghosts, so
\begin{equation}
 D_{\rm ghost}\ge-4.
\end{equation}
At $\delta=0$ this implies
\begin{equation}\label{eq:matter-bound}
 D_{\rm matter}\le4.
\end{equation}
Every nonvacuum Weyl quantum has weight at least two.  Therefore only
particle numbers $N=0,1,2$ can contribute to the centered ghost degree
selected by the Hamada polarization.  All $N\ge3$ sectors have matter weight at least six and are
contractible before any matrix-rank computation.

\section{Complete centered cohomology}\label{sec:centered}

\subsection{The relative weight-four kernel}

The complete matter Fock space at weight four is
\begin{equation}\label{eq:F4}
 \F_4=\cH_4\oplus\operatorname{Sym}^2\cH_2,
 \qquad \dim\F_4=82+55=137.
\end{equation}
The induced matter form has signature
\begin{equation}\label{eq:F4-signature}
 \operatorname{sig}(J_{\F_4})=(97,40).
\end{equation}
The relative conformal-primary conditions are
\begin{equation}\label{eq:relative-conditions}
 (D-4)|\Psi\rangle=0,
 \qquad R|\Psi\rangle=0,
 \qquad K^-_q|\Psi\rangle=0.
\end{equation}
Their exact kernel on the full space \eqref{eq:F4} is two-dimensional.  In a
spin-two magnetic basis its normalized chiral generators are
\begin{equation}\label{eq:W-square-states}
 |W_\pm^2\rangle
 =\sqrt{\frac25}|2,-2\rangle_\pm
 -\sqrt{\frac25}|1,-1\rangle_\pm
 +\frac1{\sqrt5}|0,0\rangle_\pm.
\end{equation}
They obey
\begin{equation}\label{eq:relative-I2}
 \langle W_r^2|J_{\F}|W_s^2\rangle=\delta_{rs}.
\end{equation}
This relative computation is not by itself an absolute cohomology theorem;
incoming exact states must also be excluded.

There is a simple representation-theoretic reason for the number two.  The
weight-two one-particle space is
\begin{equation}
 E_2^+\oplus E_2^-=(2,0)\oplus(0,2).
\end{equation}
The symmetric square decomposes as
\begin{equation}
 \operatorname{Sym}^2(2,0)
 \oplus\bigl[(2,0)\otimes(0,2)\bigr]
 \oplus\operatorname{Sym}^2(0,2).
\end{equation}
There is one compact scalar in each same-chirality symmetric square and no
scalar in the mixed product.  Since $E_2^\pm$ are lowest-weight spaces,
every $K^-_q$ annihilates both factors, so these two scalar singlets are
relative primaries.  By contrast, the one-particle weight-four space
contains
\begin{equation}
 (3,1)\oplus(2,1)\oplus(2,0)
\end{equation}
and its parity conjugate, with no compact scalar.  Thus representation
theory predicts the two relative candidates before the matrix calculation.
The absolute computation below is still necessary to exclude exactness and
to prove the one-particle acyclicity.  A complete Hochschild--Serre or
Kostant-style derivation of the latter ranks would be a useful conceptual
replacement for the large matrices \cite{CapBGG}.

\subsection{Absolute kernel and image}

The cutoff-complete absolute calculation supplies that missing test.  For
one chirality in the one-particle coefficient module, the centered complex
around ghost number four has exact dimensions
\begin{equation}\label{eq:one-particle-complex}
 C^3_0(290)\xrightarrow{d_3}
 C^4_0(1311)\xrightarrow{d_4}
 C^5_0(3657).
\end{equation}
The entries lie in
$\mathbb Q(i,\sqrt2,\sqrt3,\sqrt5)$.  Good-prime reduction modulo $241$,
combined with exact nilpotency, fixes the characteristic-zero ranks to
\begin{equation}
 \rank d_3=260,
 \qquad \rank d_4=1051.
\end{equation}
Since the ranks sum to $1311$,
\begin{equation}\label{eq:one-particle-zero}
 H^4_{\delta=0,N=1}=0.
\end{equation}
Parity supplies the identical result for the opposite chirality.

The lowest two-particle block begins at matter weight four.  Exterior
saturation makes its beginning
\begin{equation}\label{eq:two-particle-complex}
 0\longrightarrow C^4_0(55)
 \xrightarrow{d_4}C^5_0(385)
 \xrightarrow{d_5}C^6_0(1155).
\end{equation}
There is no incoming $C^3_0$ space: three residual ghosts have compact
degree at least $-3$ and cannot center a two-particle coefficient of minimum
weight four.  Exact nilpotency gives $d_5d_4=0$.
The two vectors \eqref{eq:W-square-states} lie in $\ker d_4$, while a
good-prime minor gives $\rank d_4\ge53$.  The displayed independent kernel
vectors give $\rank d_4\le53$.  Therefore
\begin{equation}\label{eq:two-particle-H}
 H^4_{\delta=0,N=2}
 =\operatorname{span}\{[W_+^2],[W_-^2]\}.
\end{equation}

For the vacuum coefficient module the complex is ordinary Lie-algebra
cohomology.  Since
\begin{equation}
 H^\bullet(\mathfrak{so}(4,2);\C)
 \cong\Lambda(u_3,u_5,u_7),
\end{equation}
one has
\begin{equation}\label{eq:vacuum-zero}
 H^4_{\delta=0,N=0}=0.
\end{equation}

\subsection{The canonical residual CE pairing}

The residual ghost normalization is fixed almost completely by the compact
grading.  Write
\begin{equation}
 \g=\g_{-1}\oplus\g_0\oplus\g_{+1},
 \qquad \dim(\g_{-1},\g_0,\g_{+1})=(4,7,4),
\end{equation}
and choose nonzero determinant vectors
\begin{equation}
 v_-\in\det(\g_{+1}^*),
 \qquad
 \Theta_0\in\det(\g_0^*),
 \qquad
 v_+\in\det(\g_{-1}^*),
\end{equation}
with $v_+=v_-^\dagger$.  Their product
\begin{equation}
 \Omega_\g=v_+\wedge\Theta_0\wedge v_-
\end{equation}
is the top Berezin volume on all fifteen residual ghosts.  Define the
polarized pairing by coefficient extraction,
\begin{equation}\label{eq:canonical-ghost-pairing}
 (\alpha,\beta)_{\rm gh}
 :=[\alpha^\dagger\wedge\Theta_0\wedge\beta]_{\Omega_\g}.
\end{equation}

\begin{lemma}[Canonical centered ghost pairing]\label{lem:ghost-pairing}
After fixing the orientation and the overall residual volume,
\eqref{eq:canonical-ghost-pairing} is the unique determinant-saturating
pairing compatible with the $(4,7,4)$ polarization.  It obeys
\begin{equation}\label{eq:ghost-norm}
 (v_-,v_-)_{\rm gh}=1.
\end{equation}
Because $\mathfrak{so}(4,2)$ is semisimple and hence unimodular, the
Chevalley--Eilenberg differential is graded skew for the associated
complementary-degree Poincar\'e pairing.
\end{lemma}

\begin{proof}
Each of the three determinant lines is one-dimensional.  Once their product
is oriented and normalized, coefficient extraction gives
\eqref{eq:ghost-norm} and leaves no matrix freedom.  Unimodularity makes the
top CE volume invariant, so integration by parts in the exterior algebra
gives graded skew-adjointness of the CE differential.
\end{proof}

Consequently exact cochains are orthogonal to closed complementary-degree
cochains, and the pairing descends to residual CE cohomology.

In Hamada's oscillator notation, $v_-$ is the product of the four
proper-conformal ghosts, $\Theta_0$ contains the dilation/time ghost and the
six rotation ghosts, and the bra supplies the four conjugate ghosts
\cite{HamadaBRST}.  The factor of $i$ is a real-structure convention.  The
repository exterior-algebra certificate verifies the saturation,
Hermiticity, compact degree $-4$, and unit overlap directly.

The fixed insertion is not a nondegenerate inner product on every individual
ghost-number block.  It is the polarized form of the canonical
complementary-degree CE pairing.  It is sufficient here because the incoming
two-particle space is empty and the two normalized representatives are
fixed.  Multiplying \eqref{eq:ghost-norm} by the matter restriction
\eqref{eq:relative-I2} gives their exact residual cohomological Gram matrix.
It is Hermitian and sesquilinear on the complex oscillator module; the
parity-compatible real structure gives its real restriction.  This is an
intrinsic statement about the chosen residual CE complex.  We reserve the
phrase ``transferred pure-Weyl BV cohomological pairing'' for the
cyclic-transfer statement in Section~\ref{sec:transfer}.  Showing that the
full local BV/BFV reduction induces this residual orientation and overall
volume is part of the conditional bridge.

\begin{theorem}[Complete cohomology under full residual conformal gauging]\label{thm:residual-main}
Let the coefficient module be the free bosonic Fock space
$\F(\W_+\oplus\W_-)$, let the residual algebra be
$\mathfrak{so}(4,2)$, treat $D$ as a gauge generator, and use the centered
ghost polarization \eqref{eq:ghost-vacuum-degree}.  Then the complete
minimal absolute residual cohomology at the selected ghost number four is
\begin{equation}\label{eq:main-result}
 \boxed{
 H^4_{\rm res,min}
 =\operatorname{span}\{[W_+^2],[W_-^2]\},
 \qquad G_{\rm res}=I_2.}
\end{equation}
There is no one-particle class and no higher-matter-weight class in this
minimal centered polarization.
\end{theorem}

\begin{proof}
The Cartan homotopy, Theorem~\ref{thm:cartan}, eliminates every nonzero
total compact degree.  The ghost-degree floor then reduces the centered
inventory to $N=0,1,2$.  Equations \eqref{eq:vacuum-zero},
\eqref{eq:one-particle-zero}, and \eqref{eq:two-particle-H} exhaust those
sectors.  Equations \eqref{eq:relative-I2} and \eqref{eq:ghost-norm} give the
pairing.
\end{proof}

\begin{remark}[What is positive]
The theorem does not turn the indefinite one-particle module
\eqref{eq:Jconf} into a positive particle space.  It removes the
one-particle sector from this absolute residual cohomology.  The positive
$I_2$ belongs to two ghost-dressed scalar vertex classes.
\end{remark}

\section{Descent and the dynamical/topological split}\label{sec:descent}

\subsection{Why weight four is selected}

Let $V_h$ be a scalar conformal primary of weight $h$, and let
\begin{equation}\label{eq:ghost-volume}
 \omega=\frac1{4!}\epsilon_{\mu\nu\rho\sigma}
 c^\mu c^\nu c^\rho c^\sigma
\end{equation}
be the top lowering-ghost product.  In the residual descent conventions,
\begin{equation}\label{eq:descent-transformations}
 sV_h=c^\mu\nabla_\mu V_h
 +\frac h4(\nabla_\mu c^\mu)V_h,
 \qquad
 s\omega=-\omega\nabla_\mu c^\mu,
 \qquad \omega c^\mu=0.
\end{equation}
It follows that both the integrated density and the unintegrated
ghost-dressed operator have the same closure defect,
proportional to $h/4-1$.  Hence
\begin{equation}\label{eq:descent-map}
 \boxed{[\omega V_4]\longleftrightarrow
 \left[\int_M d^4x\,\sqrt{-g}\,V_4\right].}
\end{equation}
The $-4$ compact degree of the ghost factor is therefore the conformal
counterpart of the four-dimensional volume form, rather than an accidental
normal-ordering shift.  This weight-four state--operator mechanism and a
Weyl-square example were already explicit in \cite{HamadaBRST}; the new use
here is the absolute exhaustion and pairing calculation of
Section~\ref{sec:centered}.

Three cohomological quotients appear and must not be identified without the
transfer theorem:
\begin{equation}\label{eq:three-cohomologies}
 H^4\bigl(\mathfrak{so}(4,2);\F_{\W}\bigr),
 \qquad
 H^{0,4}_{\rm loc}(s\mid d),
 \qquad
 \frac{\{\text{integrated local functionals}\}}
 {\{\text{allowed boundary functionals}\}}.
\end{equation}
The first is the residual Chevalley--Eilenberg group computed exactly in
this paper.  The second is local BRST cohomology modulo spacetime descent
\cite{BarnichBrandtHenneaux}.  The third depends on topology and boundary
conditions.  Equation \eqref{eq:descent-map} identifies representatives in
the residual construction; its promotion to an isomorphism with the strict
pure-Weyl local group is part of the conditional bridge.

\subsection{Parity basis}

Parity exchanges $W_+^2$ and $W_-^2$.  Define
\begin{equation}\label{eq:parity-basis}
 e=\frac{W_+^2+W_-^2}{\sqrt2},
 \qquad
 o=\frac{W_+^2-W_-^2}{\sqrt2}.
\end{equation}
Up to the orientation-dependent factor of $i$ in Lorentzian self-duality
conventions, descent identifies
\begin{equation}
 e\sim C_{\mu\nu\rho\sigma}C^{\mu\nu\rho\sigma},
 \qquad
 o\sim C_{\mu\nu\rho\sigma}\widetilde
 C^{\mu\nu\rho\sigma}.
\end{equation}
Because \eqref{eq:parity-basis} is orthogonal, its Gram matrix remains
$I_2$.

\subsection{Euler--Lagrange quotient}

Let
\begin{equation}
 \cE:H^4_{\rm res,min}
 \longrightarrow
 \{\text{local equation and gauge deformations}\}
\end{equation}
be the Euler--Lagrange map obtained by varying the integrated density.  In
the conventions of Section~\ref{sec:local-detour},
\begin{equation}\label{eq:EL-map}
 \cE(e)=\lambda_B B_{\mu\nu},
 \qquad \lambda_B\ne0,
 \qquad
 \cE(o)=0.
\end{equation}
The first identity is the Weyl-square variation.  Under locality,
smoothness, Lorentz invariance, and a four-derivative bound, it is the unique
nontrivial self-interaction of a single linear conformal graviton modulo
equivalence \cite{BoulangerHenneaux}.

The second identity is Chern--Weil transgression.  In a local frame,
\begin{equation}\label{eq:chern-weil}
 P_4=\Tr(R\wedge R)=dQ_3(\Gamma),
\end{equation}
where
\begin{equation}\label{eq:CS3}
 Q_3(\Gamma)=\Tr\left(
 \Gamma\wedge d\Gamma+\frac23\Gamma\wedge\Gamma\wedge\Gamma
 \right).
\end{equation}
Its variation is
\begin{equation}\label{eq:Pvariation}
 \delta P_4=2d\Tr(\delta\Gamma\wedge R),
\end{equation}
using the Bianchi identity.  On a finite cylinder, in a chosen global frame,
\begin{equation}\label{eq:finite-cylinder-P}
 \int_{[t_1,t_2]\times S^3}P_4
 =\CS[\Gamma(t_2)]-\CS[\Gamma(t_1)].
\end{equation}
Thus fixed-endpoint or compactly supported variations have no bulk
Euler--Lagrange derivative.

Equations \eqref{eq:EL-map}--\eqref{eq:finite-cylinder-P} give the exact
sequence
\begin{equation}\label{eq:EL-exact-sequence}
 0\longrightarrow\operatorname{span}\{o\}
 \longrightarrow H^4_{\rm res,min}
 \xrightarrow{\cE}\operatorname{span}\{B_{\mu\nu}\}
 \longrightarrow0.
\end{equation}
Quotienting variationally trivial local bulk functionals yields
\begin{equation}\label{eq:dynamic-quotient}
 H^4_{\rm dyn}
 =H^4_{\rm res,min}/\ker\cE
 \cong\operatorname{span}\{[C^2]\},
 \qquad
 \boxed{G_{\rm dyn,res}=I_1.}
\end{equation}
Equivalently,
\begin{equation}\label{eq:I2-split}
 \boxed{I_2=I_1^{\rm dynamical}\oplus I_1^{\rm topological}.}
\end{equation}

The odd class is only locally and variationally trivial.  A gravitational
theta parameter can affect boundaries, horizons, contact terms, large frame
transformations, and nontrivial topology \cite{FischlerKundu}.  We do not
quotient those global data without specifying the boundary problem.

\section{The local-to-residual transfer}\label{sec:transfer}

Theorem~\ref{thm:residual-main} begins with the locally reduced Weyl module.
This section states exactly what is needed to identify it with the complete
free pure-Weyl BV answer.

\subsection{The free Fock lift is formal}

\begin{proposition}[Algebraic symmetric-algebra lift]
\label{prop:Fock-lift}
Let $(C,q)$ be the algebraic one-particle free complex over characteristic
zero, with $q$ extended as a derivation to $\operatorname{Sym}C$.  Then
\begin{equation}\label{eq:Fock-lift}
 \boxed{
 H(\operatorname{Sym}C,q)
 \cong\operatorname{Sym}H(C,q).}
\end{equation}
In particular, once the one-particle coefficient module is
$\W_+\oplus\W_-$, its free many-particle cohomology is
$\F(\W_+\oplus\W_-)$.
\end{proposition}

\begin{proof}
Degree by degree, $C^{\otimes N}$ obeys the K\"unneth theorem.  Taking
symmetric coinvariants under the finite group $S_N$ is exact in
characteristic zero, so
$H(\operatorname{Sym}^NC,q)\cong\operatorname{Sym}^NH(C,q)$.  Taking the
algebraic direct sum over $N$ proves \eqref{eq:Fock-lift}.
\end{proof}

Thus the Fock-space coefficient module is not an independent field-theory
hypothesis.  The statement is algebraic; topological tensor products and
completed Fock spaces require their own continuity assumptions.

\subsection{Equivariant Cartan transfer}

Let $(C,q)$ be the gauge-fixed local free BV complex after the fifteen
conformal-Killing zero modes have been separated, and let $H=H(q)$.  Suppose
there is a strong deformation retract
\begin{equation}\label{eq:SDR}
 (H,0)\mathrel{\mathop{\rightleftarrows}^{\jmath}_{p}}(C,q),
 \qquad
 \jmath p=\id-qs-sq,
 \qquad p\jmath=\id_H.
\end{equation}
Let $\Delta$ contain the residual zero-mode differential and mixed terms, so
$Q=q+\Delta$.  With the sign convention in \eqref{eq:SDR}, homological
perturbation gives
\begin{align}\label{eq:HPL-data}
 I&=(\id+s\Delta)^{-1}\jmath,
 &P&=p(\id+\Delta s)^{-1},
 &Q_H&=p\Delta(\id+s\Delta)^{-1}\jmath.
\end{align}
If
\begin{equation}\label{eq:D-equivariance}
 [D,q]=[D,\Delta]=[D,s]=0,
 \qquad pD=D_Hp,
 \qquad D\jmath=\jmath D_H,
\end{equation}
then the transferred operators
\begin{equation}
 \iota_D^H=P\iota_DI,
 \qquad \cL_D^H=P\cL_DI
\end{equation}
obey
\begin{equation}\label{eq:transferred-Cartan}
 [Q_H,\iota_D^H]_+=\cL_D^H.
\end{equation}
Thus every transferred subspace on which $\cL_D^H$ is invertible is again
contractible.

\subsection{Cyclic transfer and the pairing}

Let $\sharp$ be the graded adjoint for the full BV/Krein pairing.  Suppose
the unperturbed retract is cyclic,
\begin{equation}\label{eq:cyclic-SDR}
 \jmath^\sharp=p,
 \qquad (s\Delta)^\sharp=\Delta s.
\end{equation}
Then taking the adjoint of \eqref{eq:HPL-data} gives
\begin{equation}
 I^\sharp=P.
\end{equation}
Since $PI=\id_H$,
\begin{equation}\label{eq:HPL-isometry}
 \boxed{I^\sharp I=\id_H.}
\end{equation}
The dressed inclusion is therefore an exact isometry: contractible local
dressing changes representatives without changing the reduced Gram matrix.

\begin{proposition}[Blockwise cyclic retract]\label{prop:cyclic-retract}
Let a finite compact-energy and $SO(4)$ block carry a nondegenerate graded
pairing, let $q^\sharp=-q$, and suppose the induced pairing on $H(q)$ is
nondegenerate.  Then the block admits a cyclic strong deformation retract.
The algebraic direct sum of these retracts may be chosen compact-degree and
$SO(4)$ equivariant.
\end{proposition}

\begin{proof}
Write $B=\im q$ and choose a nondegenerate representative space
$\cH\subset\ker q$.  In $\cH^\perp$, choose an isotropic complement $L$ to
$B$.  Then $q:L\to B$ is an isomorphism.  Define $s|_B=(q|_L)^{-1}$ and set
$s=0$ on $\cH\oplus L$.  The decomposition
$C=\cH\oplus B\oplus L$ gives $qs+sq=1-\jmath p$, while the hyperbolic
pairing of $B$ and $L$ gives the cyclic adjoint relation.  Choosing the
splitting in each finite symmetry block and taking the algebraic direct sum
preserves compact degree and $SO(4)$ type.
\end{proof}

This proposition does not provide one splitting which is simultaneously
equivariant under the noncompact generators $K^\pm$.  Such a splitting
cannot be obtained by averaging over the nonunitary conformal module, which
need not be semisimple.  Compatibility with the full $SO(4,2)$ action, the
global zero-mode split, and the BV pairing remains a field-theory
obligation.  Boundedness on a completed cylinder space is a separate
analytic question.

\subsection{Raw polynomial instantiation and measured noncompact defects}

The abstract discussion can be instantiated before choosing an analytic
completion.  In the finite polynomial conformal category, fixed total
compact energy gives the exact rational chain
\begin{equation}\label{eq:raw-polynomial-chain}
 G_n\xrightarrow{A_n}M_n\xrightarrow{B_n}E_n
 \xrightarrow{C_n}I_n,
\end{equation}
where the four slots contain the Diff$\times$Weyl ghosts, trace-free metric
and trace doublet, Bach equation/metric-antifield row, and Noether
identity/ghost-antifield row.  The Weyl scalar is written as
\begin{equation}
 \omega=\sigma+\frac14\partial\!\cdot\!\xi,
\end{equation}
so the vector arrow is the trace-free conformal-Killing operator and the
trace is an explicit doublet.  The raw matrices commute with all four
coordinate translations and special-conformal vector fields, not only with
$D\times SO(4)$.

\begin{cproposition}[Raw metric-BV retract and conformal transfer]
\label{prop:raw-BV-transfer}
In the complete energy buffer $n=2,3,4,5$, exact rational changes of basis
give maps $p_n,\jmath_n,s_n$ satisfying
\begin{equation}
 p_n\jmath_n=1,
 \qquad
 \jmath_np_n=1-q_ns_n-s_nq_n,
\end{equation}
with cohomology dimensions $10,40,82,136$.  The chosen retract is not
strictly equivariant under the noncompact generators.  For one Cartesian
component, the defect ranks are
\begin{equation}\label{eq:raw-defect-ranks}
\begin{array}{c|rrrr}
 &\rank\rho_H&\rank\delta_{\jmath}&\rank\delta_p&\rank\delta_s\\
\hline
 P_0&40&29&5&35\\
 K_0&82&80&9&97.
\end{array}
\end{equation}
Every inclusion defect is explicitly $q$-exact, every projection defect has
an explicit right $q$-homotopy, and the homotopy defect obeys the induced
projector identity.  The transferred maps obey all sixteen $[K_b,P_a]$
brackets exactly on cohomology.  Moreover,
\begin{equation}\label{eq:physical-HPL-zero}
 p\rho_a s\rho_b\jmath=0
\end{equation}
on the physical metric row for every pair of proper-conformal components;
all higher terms vanish because a second $s$ acts on the gauge row.
\end{cproposition}

Thus full equivariance of an arbitrarily chosen representative inclusion is
not needed and, for this natural exact splitting, is false.  What transfers
strictly is the conformal action on cohomology.  The calculation therefore
realizes the homotopy-equivariant outcome anticipated above rather than
silently replacing it by a strict SDR.  The exact matrices, defect
homotopies, and generated table are produced by
\path{symbolic/verify_conformal_raw_bv_transfer.py}.

\subsection{The filtered local-to-residual complex}

The possible failure modes can be displayed without compressing them into
the phrase ``no other row.''  After the conformal-Killing zero modes have
been split from the local ghosts, introduce the residual-ghost filtration
on the bicomplex
\begin{equation}\label{eq:bicomplex}
 C^{p,q}=\Lambda^p\g^*\widehat\otimes C^q_{\rm local},
 \qquad
 F^pC=\bigoplus_{r\ge p}C^{r,*}.
\end{equation}
Here $p$ is residual ghost number and $q$ is local BV degree.  A transferred
degree-one differential has the filtered expansion
\begin{equation}\label{eq:transferred-expansion}
 Q_H=q_{\rm local}+d_{\rm res}
 +\sum_{k\ge2}Q_{[k]},
\end{equation}
with bidegrees
\begin{equation}
 |q_{\rm local}|=(0,1),
 \qquad |d_{\rm res}|=(1,0),
 \qquad |Q_{[k]}|=(k,1-k).
\end{equation}
The terms $Q_{[k]}$ encode higher transferred ghost operations; assuming
that they vanish is a substantive strictness assertion, not automatic
homological bookkeeping.

The resulting spectral sequence begins
\begin{align}\label{eq:spectral-pages}
 E_1^{p,q}
 &=\Lambda^p\g^*\otimes H^q(q_{\rm local}),
 &E_2^{p,q}
 &=H^p\!\left(\g;H^q(q_{\rm local})\right),
\end{align}
and
\begin{equation}\label{eq:spectral-bidegree}
 d_r:E_r^{p,q}\longrightarrow E_r^{p+r,q-r+1}.
\end{equation}
For the desired term $E_r^{4,0}$, a higher differential can enter from
$(4-r,r-1)$ for $r=2,3,4$, and can leave toward $(4+r,1-r)$ for
$2\le r\le11$.  Thus local classes in degrees $q=1,2,3$ and antifield or
negative-degree classes down to $q=-10$ are potential modifiers.  They must
be inventoried or excluded; it is not enough to compute only the local
degree-zero row.  Since $0\le p\le15$, this filtration is bounded.  On each
finite compact-energy block the cochain spaces used here are finite, so the
spectral sequence converges.  An infinite completed state space additionally
requires a complete, separated filtration.  Even after collapse, a cyclic
splitting is needed to control extension data and the Hermitian form.

This is the standard local-BRST setting in which antifield and gauge-fixed
cohomologies must be compared carefully
\cite{BarnichBrandtHenneaux,BarnichHenneauxHurthSkenderis}.

\begin{proposition}[Single-row strictness and collapse]
\label{prop:single-row}
Suppose the global zero modes have been split exactly,
$H^q(q_{\rm local})=0$ for $q\ne0$, and the free residual perturbation
$\Delta$ has bidegree $(1,0)$.  Suppose in addition that the chosen SDR is
compatible with the local grading, so that $s$ has bidegree $(0,-1)$ and
$\jmath,p$ are supported in local degree zero.  Then
\begin{equation}\label{eq:strict-transfer}
 \boxed{Q_H=p\Delta\jmath,}
\end{equation}
all higher HPL terms vanish, and the local-to-residual spectral sequence
collapses to
\begin{equation}
 H^p\!\left(\g;H^0(q_{\rm local})\right).
\end{equation}
\end{proposition}

\begin{proof}
The homotopy $s$ lowers local degree by one.  Since both $\jmath(H)$ and the
image retained by $p$ lie in local degree zero,
\begin{equation}
 p\Delta(s\Delta)^r\jmath=0\qquad(r\ge1).
\end{equation}
The HPL series therefore reduces to \eqref{eq:strict-transfer}.  The
$E_1$ page has only the row $q=0$, so no higher differential can enter or
leave it.
\end{proof}

The proposition concerns the perturbation obtained from the \emph{free
metric BV bicomplex}.  It must not be inferred by naively restricting a
nonlinear minimal $Q$-manifold.  In the Dneprov--Grigoriev variables,
curvature-dependent terms schematically of the form $\xi\xi W$ and
$\xi\xi C$ occur in ghost transformations and can have total compact degree
zero \cite{DneprovGrigoriev}.  They belong to the nonlinear comparison and
must either be identified with the interacting deformation or removed by
the equivalence back to the free metric presentation before
Proposition~\ref{prop:single-row} is applied.

\begin{proposition}[Formal filtered-transfer criterion]\label{prop:formal-transfer}
Suppose the zero-mode split produces \eqref{eq:bicomplex}; the filtration is
convergent; $H^q(q_{\rm local})$ vanishes in every row that can enter or
leave $(4,0)$; and the transferred differential on $H^0(q_{\rm local})$ is
the strict residual Chevalley--Eilenberg differential, with all relevant
$Q_{[k]}$ absent.  Then the full cohomology in the selected filtered degree
is isomorphic, as a vector space, to the residual group
\begin{equation}
 H^4\!\left(\g;H^0(q_{\rm local})\right).
\end{equation}
If in addition the retract is compact-degree equivariant and cyclic, the
Cartan homotopy and the cohomological Hermitian pairing transfer, and the
dressed inclusion is the isometry \eqref{eq:HPL-isometry}.
\end{proposition}

\begin{proof}
Under the stated row and strictness assumptions no differential
\eqref{eq:spectral-bidegree} enters or leaves $(4,0)$, so the bounded
spectral sequence collapses there to its $E_2$ value.  The equivariant and
cyclic statements are \eqref{eq:transferred-Cartan} and
\eqref{eq:HPL-isometry}.  The cyclic splitting fixes the extension of the
pairing from the associated graded object.
\end{proof}

\subsection{End-to-end centered integration in the polynomial category}

The transferred coefficient matrices of
Proposition~\ref{prop:raw-BV-transfer} can be coupled directly to the
independently generated residual BFV/CE ghosts.  This removes the
hand-specified $E/A/L$ matrices from the final centered rank calculation.

\begin{cproposition}[Metric-to-residual integration]
\label{prop:metric-residual-integration}
In the parity-complete algebraic polynomial complex, the pipeline
\begin{equation}
 \text{metric BV rows}\longrightarrow
 H(q)\longrightarrow
 C^\bullet(\mathfrak{so}(4,2);H(q))
\end{equation}
has the exact centered dimensions and ranks
\begin{equation}\label{eq:integrated-rank-table}
\begin{array}{c|rrr|rr|r}
 &\dim C^3&\dim C^4&\dim C^5&\rank d_3&\rank d_4&\dim H^4\\
\hline
\text{vacuum}&147&407&833&116&291&0\\
\text{one particle}&580&2622&7314&520&2102&0\\
\text{two particles}&0&55&385&0&53&2.
\end{array}
\end{equation}
For the two-particle row, the next space has dimension $1155$ and
$d_5d_4=0$ exactly.  Orientation reversal acts on the two kernel vectors by
\begin{equation}
 \operatorname{diag}(-1,+1).
\end{equation}
Pulling the positive energy-two electric/magnetic curvature form back to
the metric representatives gives the raw kernel Gram matrix
\begin{equation}
 \operatorname{diag}\!\left(\frac5{64},5\right).
\end{equation}
After exact normalization and multiplication by the intrinsic unit residual
ghost norm, the representative Gram matrix is $I_2$.
\end{cproposition}

The odd and even eigenvectors are respectively the Pontryagin and
Weyl-square directions.  The calculation verifies the vector-space transfer
and positive representative form in this algebraic model.  It does not by
itself identify the curvature-coordinate normalization with the complete
cross-energy local BV/BFV symplectic pairing.  That final cyclic
normalization comparison remains visible below.  The certificate and
machine-readable ranks are produced by
\path{symbolic/verify_conformal_metric_to_residual_integration.py}.

\subsection{Two remaining field-theoretic hypotheses}

The smooth BGG theorem and Theorem~\ref{thm:algebraic-cylinder} remove the
global PDE-surjectivity and finite-mode preimage gaps.
Propositions~\ref{prop:raw-BV-transfer} and
\ref{prop:metric-residual-integration} additionally close the kernel,
image, noncompact-action, and centered-rank calculation in an explicit
polynomial metric-BV model.  Identifying that model and its representative
form with the complete field-theoretic BV/BFV complex is concentrated in
two statements.

\begin{conjecture}[Strict pure-Weyl transfer hypotheses]\label{hyp:pure-weyl}
The gauge-fixed free strict-pure-Weyl BV complex has the following
properties in the cylinder problem under consideration.
\begin{enumerate}
\item \emph{Field-theoretic BV and cyclic identification.}
The fifteen global reducibilities and their genuine dual
constraint/antifield sectors can be split without double counting.  After
nonminimal doublets are contracted, the complete nonzero-mode free local BV
complex is quasi-isomorphic, in the selected cylinder window, to the raw
polynomial detour rows of \eqref{eq:raw-polynomial-chain}, with no additional
local or antifield row entering or leaving the centered degree.  Its
cross-energy BV/Krein pairing admits a cyclic transfer whose normalization
agrees with the curvature form used in
Proposition~\ref{prop:metric-residual-integration}.
\item \emph{Residual BFV choice.}
The selected closed-universe boundary problem gauges all fifteen residual
generators and its BFV pairing induces the compact-energy four-ghost
polarization and complementary-degree determinant insertion of
Lemma~\ref{lem:ghost-pairing}, including the common normalization
\eqref{eq:ghost-norm} and the parity-compatible real structure.
\end{enumerate}
\end{conjecture}

These are not consequences of the residual rank matrices.  The raw
calculation already contains the minimal detour ghost and antifield rows and
proves homotopy-equivariant noncompact transfer; the first item asks for the
remaining identification with the complete field-theoretic BV domain and
its cross-energy cyclic form.  The second is stronger than verifying the
natural Berezin saturation inside the already chosen residual complex.

\begin{corollary}[Conditional algebraic strict-pure-Weyl result]\label{cor:conditional}
If Field-Theoretic Hypothesis~\ref{hyp:pure-weyl} holds, the complete free
classical local-plus-residual BV cohomology in the algebraic cylinder
category and selected degree is \eqref{eq:main-result}.  Its transferred
cohomological Hermitian form is $I_2$, and its local Euler--Lagrange quotient
is \eqref{eq:dynamic-quotient} with form $I_1$.
\end{corollary}

\begin{proof}
Theorems~\ref{thm:curvature-isomorphism} and
\ref{thm:algebraic-cylinder} identify the one-particle coefficient module.
Proposition~\ref{prop:Fock-lift} gives its free Fock lift.
Propositions~\ref{prop:raw-BV-transfer} and
\ref{prop:metric-residual-integration} give the strict induced residual
action and centered cohomology in the algebraic polynomial model.  Item 1
identifies this model and its cyclic form with the complete local BV
complex.  Item 2 fixes the selected ghost pairing and normalization.
Equation~\eqref{eq:HPL-isometry} then identifies the residual representative
Gram matrix with the transferred pure-Weyl BV cohomological form.
\end{proof}

The global smooth BGG theorem, curvature isomorphism
\eqref{eq:curvature-isomorphism}, action factorization
\eqref{eq:Bfactor}, raw metric-BV transfer, and integrated centered ranks are
exact inputs.  The two hypotheses above state the remaining obligations
without folding the cross-energy cyclic pairing or BFV boundary choice into
the metric-to-curvature theorem.  A Hilbert, Krein, distributional, or
completed Fock statement
would require separate continuity, closed-range, and convergence results.

\section{Scope, interpretation, and open questions}\label{sec:scope}

\paragraph{Residual gauge choice.}
The result is specific to quotienting the full residual $SO(4,2)$ action on
the compact cylinder.  If some generators are retained as global or
asymptotic charges, the physical complex is different.  In particular,
Cartan contraction cannot be used to erase nonzero $D$ energy when $D$ is
the physical Hamiltonian.

\paragraph{No positive graviton Fock claim.}
The unreduced module remains indefinite and the one-particle residual
cohomology is zero.  The positive classes are weight-four scalar vertices.
Calling $I_2$ a positive graviton Hilbert space would therefore misidentify
both the cohomological degree and the state--operator map.

\paragraph{Classical and free.}
The result concerns the free classical complex and its finite residual
reduction.  It does not show that the interaction preserves the reduced
space or pairing.  Nor does it address the quantum master equation.  A local
Diff$\times$Weyl anomaly could obstruct quantum nilpotency even if
Corollary~\ref{cor:conditional} is fully established.

\paragraph{Vertex versus dynamics.}
The two residual classes do not define two independent local gravitational
dynamics.  The parity-even class is the Weyl action direction; the
parity-odd class is a theta direction, locally related by a canonical
boundary transformation.  The positive dynamical quotient is therefore
one-dimensional.

\paragraph{The remaining classical bridge.}
The all-level smooth metric-to-$\mathcal W_{\rm geom}$ identification and
the algebraic $E/A/L$ curvature intertwiner are exact.  The raw polynomial
metric-BV complex now supplies the antifield-row inventory, explicit
noncompact homotopies, strict induced conformal action, and end-to-end
centered ranks.  Two sharper obligations remain: identify its
curvature-coordinate form with the complete cross-energy cyclic local-BV
pairing, and derive the chosen four-ghost polarization and normalized
pairing from the residual BFV boundary problem.  The direct Taub argument
does not require first identifying a degree-three BGG representative,
although that geometric comparison remains valuable.  Extension to a
chosen Hilbert or Krein completion is a separate analytic question.  The
exact ranks in Table~\ref{tab:detour-ranks} remain independent convention and
regression checks.

\section{Conclusion}\label{sec:conclusion}

The conventional free oscillator description of pure conformal gravity on
$\R\times S^3$ is indefinite.  Residual conformal reduction changes the
question rather than re-signing those oscillators.  The absolute residual
complex carries a Cartan contraction which eliminates every noncentered
compact degree.  At the ghost number four selected by the centered
polarization this turns the apparent
infinite problem into the three sectors $N=0,1,2$.  The first two are
acyclic; the last has exactly two classes and positive Gram matrix:
\begin{equation*}
 H^4_{\rm res,min}
 =\operatorname{span}\{[W_+^2],[W_-^2]\},
 \qquad G_{\rm res}=I_2.
\end{equation*}
Descent identifies them as the Weyl-square and Pontryagin vertex directions,
and the quotient by variationally trivial bulk functionals leaves one
positive dynamical direction.

The result is both stronger and narrower than a sampled positivity
calculation.  It is exhaustive for the stated minimal residual complex, and
the same centered ranks now arise from an independent raw polynomial
metric-BV transfer, but it concerns vertices rather than propagating
particles.  Its promotion to the complete field-theoretic pure-Weyl BV/BFV
theory, including the cyclic normalization and boundary choice, is
controlled by the explicit hypotheses of
Corollary~\ref{cor:conditional}.  Quantum consistency and interaction
descent are separate questions.

\appendix

\section{Conventions}\label{app:conventions}

The cylinder radius is one and the Lorentzian signature is $(-,+,+,+)$.
Compact energy is the eigenvalue of $D=i\partial_t$ in the oscillator
convention.  The compact algebra is
$SO(4)\cong SU(2)_L\times SU(2)_R$.  A label $(j_L,j_R)$ has dimension
$(2j_L+1)(2j_R+1)$.  Positive chirality means the lowest Weyl primary
$(2;2,0)$; parity exchanges $j_L$ and $j_R$.

The normalized module convention is \eqref{eq:Jconf}.  The reduced curvature
action \eqref{eq:Sred} differs by an overall factor and sign from the common
$-\int C^2$ oscillator convention.  All moment-map comparisons derive the
symplectic form and Taub kernel from the same action, so the overall action
factor cancels.  Lorentzian self-duality may insert a factor of $i$ in the
odd parity combination; no conclusion depends on that convention.

The residual generator degrees are $(0^7,-1^4,+1^4)$ and dual ghost degrees
are their negatives.  Ghost number and compact degree are independent.
The selected absolute residual degree in this paper is ghost number four,
fixed by \eqref{eq:ghost-vacuum-degree}; it is not ordinary ghost number
zero.  We call it a candidate physical degree only under the full residual
gauging specified in Section~\ref{sec:D-gauge}.

\section{Exact rank certification}\label{app:ranks}

The one-particle matrices contain radicals.  They are represented over
\begin{equation}
 K=\mathbb Q(i,\sqrt2,\sqrt3,\sqrt5).
\end{equation}
For a lower bound on rank, the calculation maps this field to
$\mathbb F_{241}$ by
\begin{equation}
 i\mapsto64,
 \qquad \sqrt2\mapsto22,
 \qquad \sqrt3\mapsto56,
 \qquad \sqrt5\mapsto103.
\end{equation}
Each image obeys its defining square relation.  A nonzero minor after this
good-prime reduction is a nonzero characteristic-zero minor.  Exact
nilpotency provides the complementary upper bound
\begin{equation}
 \rank d_{q-1}+\rank d_q\le\dim C^q.
\end{equation}
For the one-particle complex the modular lower bounds saturate this inequality.
For the two-particle complex the two explicit independent kernel vectors
provide the upper bound which the modular rank saturates.  Thus the ranks
quoted in Section~\ref{sec:centered} are exact characteristic-zero ranks,
not probabilistic modular estimates.

\section{Claim-status and verification index}\label{app:verification}

The residual computation used in the first draft was frozen at repository
commit
\begin{center}
\texttt{e928f257c25099eb534eb34109dfc1dc6a3127a1}.
\end{center}
The global BGG bridge was added in the subsequent theorem revision based at
\begin{center}
\texttt{c471b99f5e3708e692b1c25238f6272c9e29b48f}.
\end{center}
The companion snapshot
records both layers, the environment, hashes, expected-failure guards, and
command lines.  The complete paper verification is run with
\begin{verbatim}
python3 symbolic/verify_conformal_paper_free.py
\end{verbatim}

\begin{center}
\footnotesize
\begin{tabular}{>{\raggedright\arraybackslash}p{0.25\textwidth}
                >{\raggedright\arraybackslash}p{0.17\textwidth}
                >{\raggedright\arraybackslash}p{0.48\textwidth}}
\toprule
Claim & Status & Evidence / boundary\\
\midrule
Formal pure-Weyl detour identities & Exact & Branson--Gover theorem plus
action-normalized factorization and finite scalar-block audits.\\
Global smooth deformation cohomology and curvature isomorphism & Exact
smooth theorem & The flat BGG fine resolution and
$H^\bullet(\R\times S^3)$ give $H^1_{\rm def}=H^2_{\rm def}=0$ and hence
\eqref{eq:curvature-isomorphism}.  This does not establish a completed
Hilbert/Krein isomorphism or the BV/BFV zero-mode split.\\
Euclidean polynomial detour ranks, degrees $2$--$6$ & Exact finite
certificate & Rational matrices verify kernels, ranks, and the $15$ residual
zero modes; they are an independent convention audit rather than the source
of the global theorem.\\
Algebraic cylinder realization and stable generator coefficients & Exact
symbolic theorem & Full coordinate Weyl/Bach calculation at symbolic $n$,
nonzero $E/A/L$ pivots, explicit same-block metric preimages, character
exhaustion, invariant form, and cutoff-stable conformal brackets.\\
Residual Cartan homotopy reduction & Theorem & Explicit homotopy
$\iota_D/\delta$ on every nonzero total degree; this is a quotient only in
the full residual-gauging problem.\\
Centered vacuum and one-particle sectors & Exact & Semisimple Lie-algebra
cohomology for the vacuum; cutoff-complete absolute rank certificate for
one particle.\\
Centered two-particle sector and $G_{\rm res}=I_2$ & Exact residual theorem &
Independent relative and absolute calculations; empty incoming absolute
space and exact rank $53$.\\
Vertex descent and dynamical/topological split & Exact locally & Weight-four
descent, Chern--Weil transgression, and Euler--Lagrange quotient.  Global
theta effects are retained.\\
Full local-plus-residual pure-Weyl BV result & Conditional & Requires the
two parts of Field-Theoretic Hypothesis~\ref{hyp:pure-weyl}: full
equivariant cyclic BV transfer with row control, and the selected residual
BFV polarization and normalization.\\
Interpretation as a physical state space & Choice-dependent & Requires a
boundary and gauge principle that treats all fifteen generators, including
$D$, as gauge rather than global charges.\\
Quantum anomaly freedom and interacting probability & Not claimed & Outside
the scope of this paper.\\
\bottomrule
\end{tabular}
\end{center}

\begin{thebibliography}{99}

\bibitem{BransonGover}
T.~Branson and A.~R.~Gover,
``The conformal deformation detour complex for the obstruction tensor,''
\href{https://arxiv.org/abs/math/0605192}{arXiv:math/0605192}.

\bibitem{GoverPeterson}
A.~R.~Gover and L.~J.~Peterson,
``The ambient obstruction tensor and the conformal deformation complex,''
Pacific J. Math. \textbf{226} (2006) 309--351,
\href{https://arxiv.org/abs/math/0408229}{arXiv:math/0408229}.

\bibitem{GoverSombergSoucek}
A.~R.~Gover, P.~Somberg, and V.~Sou\v{c}ek,
``Yang--Mills detour complexes and conformal geometry,''
Commun. Math. Phys. \textbf{278} (2008) 307--327,
\href{https://arxiv.org/abs/math/0606401}{arXiv:math/0606401}.

\bibitem{GreenHyperbolic}
M.~Benini, G.~Musante, and A.~Schenkel,
``Green hyperbolic complexes on Lorentzian manifolds,''
Commun. Math. Phys. (2023),
\href{https://doi.org/10.1007/s00220-023-04807-5}{doi:10.1007/s00220-023-04807-5},
\href{https://arxiv.org/abs/2207.04069}{arXiv:2207.04069}.

\bibitem{CapBGG}
A.~\v{C}ap,
``Overdetermined systems, conformal geometry, and the BGG complex,''
in \emph{Symmetries and Overdetermined Systems of Partial Differential
Equations}, IMA Vol. Math. Appl. \textbf{144} (Springer, 2008) 1--24,
\href{https://arxiv.org/abs/math/0610225}{arXiv:math/0610225}.

\bibitem{CapDeformations}
A.~\v{C}ap,
``Infinitesimal automorphisms and deformations of parabolic geometries,''
J. Eur. Math. Soc. \textbf{10} (2008) 415--437,
\href{https://arxiv.org/abs/math/0508535}{arXiv:math/0508535}.

\bibitem{CapHuPoincare}
A.~\v{C}ap and K.~Hu,
``Bounded Poincar\'e operators for twisted and BGG complexes,''
J. Math. Pures Appl. \textbf{179} (2023) 253--276,
\href{https://arxiv.org/abs/2304.07185}{arXiv:2304.07185}.

\bibitem{HamadaHorata}
K.-j.~Hamada and S.~Horata,
``Conformal algebra and physical states in non-critical 3-brane on
$\R\times S^3$,''
Prog. Theor. Phys. \textbf{110} (2003) 1169--1210,
\href{https://arxiv.org/abs/hep-th/0307008}{arXiv:hep-th/0307008}.

\bibitem{HamadaBuildingBlocks}
K.-j.~Hamada,
``Building blocks of physical states in a non-critical 3-brane on
$\R\times S^3$,''
Int. J. Mod. Phys. A \textbf{20} (2005) 5353--5398,
\href{https://arxiv.org/abs/hep-th/0402136}{arXiv:hep-th/0402136}.

\bibitem{HamadaCFT}
K.-j.~Hamada,
``Conformal field theory on $\R\times S^3$ from quantized gravity,''
Int. J. Mod. Phys. A \textbf{24} (2009) 3073--3110,
\href{https://arxiv.org/abs/0811.1647}{arXiv:0811.1647}.

\bibitem{HamadaBRST}
K.-j.~Hamada,
``BRST analysis of physical fields and states for 4D quantum gravity on
$\R\times S^3$,''
Phys. Rev. D \textbf{86} (2012) 124006,
\href{https://arxiv.org/abs/1202.4538}{arXiv:1202.4538}.

\bibitem{KuzenkoPonds}
S.~M.~Kuzenko and M.~Ponds,
``Conformal geometry and (super)conformal higher-spin gauge theories,''
JHEP \textbf{05} (2019) 113,
\href{https://arxiv.org/abs/1902.08010}{arXiv:1902.08010}.

\bibitem{BeccariaBekaertTseytlin}
M.~Beccaria, X.~Bekaert, and A.~A.~Tseytlin,
``Partition function of free conformal higher spin theory,''
JHEP \textbf{08} (2014) 113,
\href{https://arxiv.org/abs/1406.3542}{arXiv:1406.3542}.

\bibitem{BoulangerHenneaux}
N.~Boulanger and M.~Henneaux,
``A derivation of Weyl gravity,''
Annalen Phys. \textbf{10} (2001) 935--964,
\href{https://arxiv.org/abs/hep-th/0106065}{arXiv:hep-th/0106065}.

\bibitem{BarnichBrandtHenneaux}
G.~Barnich, F.~Brandt, and M.~Henneaux,
``Local BRST cohomology in gauge theories,''
Phys. Rept. \textbf{338} (2000) 439--569,
\href{https://arxiv.org/abs/hep-th/0002245}{arXiv:hep-th/0002245}.

\bibitem{BarnichAntifield}
G.~Barnich, F.~Brandt, and M.~Henneaux,
``Local BRST cohomology in the antifield formalism: I. General theorems,''
Commun. Math. Phys. \textbf{174} (1995) 57--92,
\href{https://arxiv.org/abs/hep-th/9405109}{arXiv:hep-th/9405109}.

\bibitem{BarnichHenneauxHurthSkenderis}
G.~Barnich, M.~Henneaux, T.~Hurth, and K.~Skenderis,
``Cohomological analysis of gauged-fixed gauge theories,''
Phys. Lett. B \textbf{492} (2000) 376--384,
\href{https://arxiv.org/abs/hep-th/9910201}{arXiv:hep-th/9910201}.

\bibitem{MetsaevBRSTBV}
R.~R.~Metsaev,
``BRST--BV approach to conformal fields,''
J. Phys. A \textbf{49} (2016) 175401,
\href{https://arxiv.org/abs/1511.01836}{arXiv:1511.01836}.

\bibitem{DneprovGrigoriev}
I.~Dneprov and M.~Grigoriev,
``Presymplectic BV--AKSZ formulation of conformal gravity,''
\href{https://doi.org/10.1140/epjc/s10052-022-11082-6}
{doi:10.1140/epjc/s10052-022-11082-6},
\href{https://arxiv.org/abs/2208.02933}{arXiv:2208.02933}.

\bibitem{BarnichBekaertGrigoriev}
G.~Barnich, X.~Bekaert, and M.~Grigoriev,
``Notes on conformal invariance of gauge fields,''
J. Phys. A \textbf{48} (2015) 505402,
\href{https://arxiv.org/abs/1506.00595}{arXiv:1506.00595}.

\bibitem{AltasTekinLinearization}
E.~Altas and B.~Tekin,
``Linearization instability for generic gravity in AdS,''
Phys. Rev. D \textbf{97} (2018) 024028,
\href{https://arxiv.org/abs/1705.10234}{arXiv:1705.10234}.

\bibitem{AltasTekinTaub}
E.~Altas and B.~Tekin,
``Second order perturbation theory in general relativity: Taub charges as
integral constraints,''
Phys. Rev. D \textbf{99} (2019) 104078,
\href{https://arxiv.org/abs/1903.11982}{arXiv:1903.11982}.

\bibitem{ChevalleyEilenberg}
C.~Chevalley and S.~Eilenberg,
``Cohomology theory of Lie groups and Lie algebras,''
Trans. Amer. Math. Soc. \textbf{63} (1948) 85--124.

\bibitem{FischlerKundu}
W.~Fischler and S.~Kundu,
``Physical effects of the gravitational $\Theta$-parameter,''
JHEP \textbf{05} (2017) 098,
\href{https://arxiv.org/abs/1612.06010}{arXiv:1612.06010}.

\end{thebibliography}

\end{document}
